% TOP_PROBLEM: {"id": 3078, "title": "Simplexity of the n-cube", "category_id": 6, "category": {"id": 6, "name": "geometry", "display_name": "Geometry", "description": "Euclidean and non-Euclidean geometry, geometric structures.", "slug": "geometry", "order_index": 6, "created_at": "2026-07-31T15:26:25.671Z"}, "status": "open", "problem_number": "OPG-1820", "set_id": 12, "statusQualification": "T(n) permits additional vertices and non-face-to-face intersections; the stricter vertex triangulation number T*(n) is distinguished explicitly."}
% FIRST_POSTED: 2026-10-01
% ENTRY_KIND: statement_only
% UnsolvedMath ID 3078; source catalogue code OPG-1820
% !TeX program = lualatex
% Definition and sources only; this file is not an LLM solution attempt.
\documentclass[11pt,a4paper]{article}
\usepackage[margin=25mm]{geometry}
\usepackage{fontspec}
\setmainfont{lmroman10-regular.otf}[BoldFont=lmroman10-bold.otf,
 ItalicFont=lmroman10-italic.otf,BoldItalicFont=lmroman10-bolditalic.otf]
\usepackage{amsmath,amssymb,mathtools}
\usepackage{unicode-math}
\setmathfont{latinmodern-math.otf}
\AtBeginDocument{\let\setminus\smallsetminus}
\usepackage{xurl,hyperref}
\hypersetup{colorlinks=true,urlcolor=blue}
\setcounter{secnumdepth}{0}
\setlength{\parindent}{0pt}
\setlength{\parskip}{5pt}
\setlength{\emergencystretch}{3em}
\begin{document}
\section*{Simplexity of the n-cube}\label{3078}
{\small UnsolvedMath ID 3078 · OPG-1820 · Geometry · OpenGarden}
\subsection{Definitions and mathematical statement}
Let $n\ge1$ and let $Q_n=[0,1]^n\subset\mathbb R^n$.
An $n$-simplex is the convex hull of $n+1$ affinely independent points.
A decomposition of $Q_n$ means a finite family
$\mathcal D=\{\Delta_1,\ldots,\Delta_m\}$ of closed $n$-simplices
contained in $Q_n$, with union $Q_n$ and pairwise disjoint interiors.
Define
\[
 T(n)=\min\{m:Q_n\text{ has such a decomposition into }m\text{ simplices}\}.
\]
Determine $T(n)$ as a function of dimension, by matching constructive
upper bounds with lower bounds valid for every allowed decomposition.

The convention here follows the source's explicit definition of $T$:
vertices of the simplices may be arbitrary points of the cube, and
intersections of distinct simplices need not be entire common faces.
For comparison, its separate parameter $T^*(n)$ imposes both additional
requirements that every simplex vertex be a cube vertex and that the
simplices form a face-to-face triangulation. Thus $T(n)\le T^*(n)$, and
these quantities must not be silently identified.

For any permutation $\sigma$ of $\{1,\ldots,n\}$, the region
$0\le x_{\sigma(1)}\le\cdots\le x_{\sigma(n)}\le1$ is an $n$-simplex.
The resulting $n!$ simplices give a decomposition, ensuring the minimization
is over a nonempty class. Scaling or translating the cube does not
alter either minimum. The target counts full-dimensional pieces; lower-dimensional
faces and their subdivisions are not counted as additional simplices.
\subsection{Short English statement}
What is the minimum number of full-dimensional simplices needed to decompose a cube, under the source’s unrestricted decomposition convention?
\subsection{Sources}
\begin{itemize}
\item[S1] ULAM.AI and the UnsolvedMath Contributors, supplied problems.json, record 3078 (OPG-1820), OpenGarden collection. Catalogue identity and source transcription; CC BY 4.0.
\url{https://huggingface.co/datasets/ulamai/UnsolvedMath}.
\item[S2] Open Problem Garden, Simplexity of the n-cube. Original problem and discussion; the mathematical formulation below is expanded from the supplied source record.
\url{http://www.openproblemgarden.org/op/simplexity_of_the_cube}.
\end{itemize}
\end{document}
